% year: תשע"ח
% type: מבחן
% semester: א
% moed: ב
% number: 1א
% points: 10
% categories: sequences
% source: exams/מבחנים/תשעח/מועד ב תשעח.pdf

\begin{question}
(10 נק') חשבו את הגבול הבא:
גבול הסדרה $\displaystyle\lim_{n\to\infty}\left(\sqrt{(n^2+2)(n^2-4)}-\sqrt{n^4-7}\right)$.
\end{question}

\begin{hint}
הכפילו וחלקו בצמוד $\sqrt{(n^2+2)(n^2-4)}+\sqrt{n^4-7}$.
\end{hint}

\begin{solution}
לכל $n\ge2$ שני השורשים מוגדרים. נכפיל ונחלק בצמוד:
\[ \sqrt{(n^2+2)(n^2-4)}-\sqrt{n^4-7}=\frac{(n^2+2)(n^2-4)-(n^4-7)}{\sqrt{(n^2+2)(n^2-4)}+\sqrt{n^4-7}}. \]
במונה: $(n^2+2)(n^2-4)=n^4-2n^2-8$, ולכן המונה הוא $n^4-2n^2-8-n^4+7=-2n^2-1$. נחלק מונה ומכנה ב-$n^2$:
\[ \frac{-2n^2-1}{\sqrt{(n^2+2)(n^2-4)}+\sqrt{n^4-7}}=\frac{-2-\frac1{n^2}}{\sqrt{\left(1+\frac2{n^2}\right)\left(1-\frac4{n^2}\right)}+\sqrt{1-\frac7{n^4}}}. \]
מאריתמטיקה של גבולות ורציפות השורש, המונה שואף ל-$-2$ והמכנה ל-$1+1=2$. לכן
\[ \boxed{\lim_{n\to\infty}\left(\sqrt{(n^2+2)(n^2-4)}-\sqrt{n^4-7}\right)=-1} \]
\end{solution}
